% Options for packages loaded elsewhere
\PassOptionsToPackage{unicode}{hyperref}
\PassOptionsToPackage{hyphens}{url}
%
\documentclass[
  brazilian,
]{article}
\usepackage{lmodern}
\usepackage{amssymb,amsmath}
\usepackage{ifxetex,ifluatex}
\ifnum 0\ifxetex 1\fi\ifluatex 1\fi=0 % if pdftex
  \usepackage[T1]{fontenc}
  \usepackage[utf8]{inputenc}
  \usepackage{textcomp} % provide euro and other symbols
\else % if luatex or xetex
  \usepackage{unicode-math}
  \defaultfontfeatures{Scale=MatchLowercase}
  \defaultfontfeatures[\rmfamily]{Ligatures=TeX,Scale=1}
\fi
% Use upquote if available, for straight quotes in verbatim environments
\IfFileExists{upquote.sty}{\usepackage{upquote}}{}
\IfFileExists{microtype.sty}{% use microtype if available
  \usepackage[]{microtype}
  \UseMicrotypeSet[protrusion]{basicmath} % disable protrusion for tt fonts
}{}
\makeatletter
\@ifundefined{KOMAClassName}{% if non-KOMA class
  \IfFileExists{parskip.sty}{%
    \usepackage{parskip}
  }{% else
    \setlength{\parindent}{0pt}
    \setlength{\parskip}{6pt plus 2pt minus 1pt}}
}{% if KOMA class
  \KOMAoptions{parskip=half}}
\makeatother
\usepackage{xcolor}
\IfFileExists{xurl.sty}{\usepackage{xurl}}{} % add URL line breaks if available
\IfFileExists{bookmark.sty}{\usepackage{bookmark}}{\usepackage{hyperref}}
\hypersetup{
  pdftitle={Parte III: Processamento de Linguagem Natural para Linguistas},
  pdfauthor={CiberExt 26-29 · FEELT38103 · Universidade Federal de Uberlândia},
  pdflang={pt-BR},
  hidelinks,
  pdfcreator={LaTeX via pandoc}}
\urlstyle{same} % disable monospaced font for URLs
\usepackage{color}
\usepackage{fancyvrb}
\newcommand{\VerbBar}{|}
\newcommand{\VERB}{\Verb[commandchars=\\\{\}]}
\DefineVerbatimEnvironment{Highlighting}{Verbatim}{commandchars=\\\{\}}
% Add ',fontsize=\small' for more characters per line
\newenvironment{Shaded}{}{}
\newcommand{\AlertTok}[1]{\textcolor[rgb]{1.00,0.00,0.00}{\textbf{#1}}}
\newcommand{\AnnotationTok}[1]{\textcolor[rgb]{0.38,0.63,0.69}{\textbf{\textit{#1}}}}
\newcommand{\AttributeTok}[1]{\textcolor[rgb]{0.49,0.56,0.16}{#1}}
\newcommand{\BaseNTok}[1]{\textcolor[rgb]{0.25,0.63,0.44}{#1}}
\newcommand{\BuiltInTok}[1]{#1}
\newcommand{\CharTok}[1]{\textcolor[rgb]{0.25,0.44,0.63}{#1}}
\newcommand{\CommentTok}[1]{\textcolor[rgb]{0.38,0.63,0.69}{\textit{#1}}}
\newcommand{\CommentVarTok}[1]{\textcolor[rgb]{0.38,0.63,0.69}{\textbf{\textit{#1}}}}
\newcommand{\ConstantTok}[1]{\textcolor[rgb]{0.53,0.00,0.00}{#1}}
\newcommand{\ControlFlowTok}[1]{\textcolor[rgb]{0.00,0.44,0.13}{\textbf{#1}}}
\newcommand{\DataTypeTok}[1]{\textcolor[rgb]{0.56,0.13,0.00}{#1}}
\newcommand{\DecValTok}[1]{\textcolor[rgb]{0.25,0.63,0.44}{#1}}
\newcommand{\DocumentationTok}[1]{\textcolor[rgb]{0.73,0.13,0.13}{\textit{#1}}}
\newcommand{\ErrorTok}[1]{\textcolor[rgb]{1.00,0.00,0.00}{\textbf{#1}}}
\newcommand{\ExtensionTok}[1]{#1}
\newcommand{\FloatTok}[1]{\textcolor[rgb]{0.25,0.63,0.44}{#1}}
\newcommand{\FunctionTok}[1]{\textcolor[rgb]{0.02,0.16,0.49}{#1}}
\newcommand{\ImportTok}[1]{#1}
\newcommand{\InformationTok}[1]{\textcolor[rgb]{0.38,0.63,0.69}{\textbf{\textit{#1}}}}
\newcommand{\KeywordTok}[1]{\textcolor[rgb]{0.00,0.44,0.13}{\textbf{#1}}}
\newcommand{\NormalTok}[1]{#1}
\newcommand{\OperatorTok}[1]{\textcolor[rgb]{0.40,0.40,0.40}{#1}}
\newcommand{\OtherTok}[1]{\textcolor[rgb]{0.00,0.44,0.13}{#1}}
\newcommand{\PreprocessorTok}[1]{\textcolor[rgb]{0.74,0.48,0.00}{#1}}
\newcommand{\RegionMarkerTok}[1]{#1}
\newcommand{\SpecialCharTok}[1]{\textcolor[rgb]{0.25,0.44,0.63}{#1}}
\newcommand{\SpecialStringTok}[1]{\textcolor[rgb]{0.73,0.40,0.53}{#1}}
\newcommand{\StringTok}[1]{\textcolor[rgb]{0.25,0.44,0.63}{#1}}
\newcommand{\VariableTok}[1]{\textcolor[rgb]{0.10,0.09,0.49}{#1}}
\newcommand{\VerbatimStringTok}[1]{\textcolor[rgb]{0.25,0.44,0.63}{#1}}
\newcommand{\WarningTok}[1]{\textcolor[rgb]{0.38,0.63,0.69}{\textbf{\textit{#1}}}}
\setlength{\emergencystretch}{3em} % prevent overfull lines
\providecommand{\tightlist}{%
  \setlength{\itemsep}{0pt}\setlength{\parskip}{0pt}}
\setcounter{secnumdepth}{-\maxdimen} % remove section numbering
\ifxetex
  % Load polyglossia as late as possible: uses bidi with RTL langages (e.g. Hebrew, Arabic)
  \usepackage{polyglossia}
  \setmainlanguage[]{brazilian}
\else
  \usepackage[shorthands=off,main=brazilian]{babel}
\fi

\title{Parte III: Processamento de Linguagem Natural para Linguistas}
\usepackage{etoolbox}
\makeatletter
\providecommand{\subtitle}[1]{% add subtitle to \maketitle
  \apptocmd{\@title}{\par {\large #1 \par}}{}{}
}
\makeatother
\subtitle{Unidade 1 - Respostas dos Exercícios}
\author{CiberExt 26-29 · FEELT38103 · Universidade Federal de
Uberlândia}
\date{Agosto de 2026}

\begin{document}
\maketitle

\hypertarget{gabarito---muxf3dulo-1-processando-diversos-idiomas}{%
\subsection{Gabarito - Módulo 1: Processando Diversos
Idiomas}\label{gabarito---muxf3dulo-1-processando-diversos-idiomas}}

\textbf{Resposta 1:} Os dois principais desafios estruturais são: 1.
\textbf{Morfologia complexa:} Idiomas aglutinativos (como finlandês e
turco) podem combinar a raiz, o tempo verbal, marcadores de caso e
pronomes em uma única palavra, diferentemente do inglês. 2.
\textbf{Ordem das palavras (Sintaxe):} Enquanto o inglês segue uma ordem
estrita Sujeito-Verbo-Objeto (SVO), muitos outros idiomas possuem ordens
flexíveis ou adotam estruturas diferentes, como Sujeito-Objeto-Verbo
(SOV), comum no japonês.

\textbf{Resposta 2:}

\begin{Shaded}
\begin{Highlighting}[]
\ImportTok{import}\NormalTok{ stanza}

\CommentTok{\# Baixa o modelo de linguagem para o finlandês}
\NormalTok{stanza.download(lang}\OperatorTok{=}\StringTok{\textquotesingle{}fi\textquotesingle{}}\NormalTok{)}

\CommentTok{\# Inicializa o pipeline apenas com o tokenizador }
\CommentTok{\# e o etiquetador de classes gramaticais}
\NormalTok{nlp\_fi }\OperatorTok{=}\NormalTok{ stanza.Pipeline(lang}\OperatorTok{=}\StringTok{\textquotesingle{}fi\textquotesingle{}}\NormalTok{, processors}\OperatorTok{=}\StringTok{\textquotesingle{}tokenize, pos\textquotesingle{}}\NormalTok{)}
\end{Highlighting}
\end{Shaded}

\textbf{Resposta 3:} Deve-se utilizar o método \texttt{to\_dict()}. Este
método converte o objeto \emph{Sentence} em uma lista de dicionários
onde cada dicionário representa um único objeto \emph{Token}, contendo
as anotações linguísticas na forma de pares chave-valor (ex: `id',
`text', `lemma', `upos').

\textbf{Resposta 4:} A maneira mais eficiente é não passar string por
string, mas coletá-las em uma lista e pré-empacotá-las em objetos
\emph{Document}. O passo a passo seria: 1. Reunir os textos como strings
em uma lista padrão (ex: \texttt{str\_docs}). 2. Utilizar uma
compreensão de lista para criar uma nova lista instanciando documentos
vazios preenchidos apenas com o texto:
\texttt{docs\_in\ =\ {[}stanza.Document({[}{]},\ text=doc)\ for\ doc\ in\ str\_docs{]}}.
3. Passar a lista completa \texttt{docs\_in} diretamente para o objeto
do modelo de linguagem (pipeline) de uma só vez para realizar a
anotação.

\textbf{Resposta 5:} A limitação é que \textbf{o idioma em questão deve
ser suportado nativamente tanto pelo Stanza quanto pelo spaCy}. Se
tentarmos processar o idioma Wolof, por exemplo, não será possível usar
o \texttt{spacy-stanza}, pois o spaCy não possui suporte a esse idioma
em seu framework de base.

\begin{center}\rule{0.5\linewidth}{0.5pt}\end{center}

\end{document}
